\documentclass[10pt]{article}
\usepackage[verbose, a4paper, hmargin=2.5cm, vmargin=2.5cm]{geometry}

\usepackage{fontspec}
\usepackage{ctex}
\usepackage{paratype}

\usepackage{amsmath}
\usepackage{amssymb}
\usepackage{amsfonts}
\usepackage{esint}
\usepackage{stmaryrd}

\usepackage[mathbf=sym]{unicode-math}
\setmainfont{STIX Two Text}
\setmathfont{STIX Two Math}

\setsansfont{Arial}
\setmonofont{Courier New}

\usepackage{makecell}
\usepackage{wasysym}
\usepackage{pifont}
\usepackage[export]{adjustbox}
\usepackage{graphicx}
\usepackage{mdframed}
\usepackage{booktabs,array,multirow}
\usepackage{tabularx}
\usepackage{hyperref}
\hypersetup{colorlinks=true, linkcolor=blue, filecolor=magenta, urlcolor=cyan,}
\urlstyle{same}
\usepackage[most]{tcolorbox}
\definecolor{mygray}{RGB}{240,240,240}
\tcbset{
  colback=mygray,
  boxrule=0pt,
}
\graphicspath{ {./images/} }
\newcommand{\HRule}{\begin{center}\rule{0.9\linewidth}{0.2mm}\end{center}}
\newcommand{\customfootnote}[1]{
  \let\thefootnote\relax\footnotetext{#1}
}
\newcommand{\lt}{\mathrel{<}}
\newcommand{\gt}{\mathrel{>}}
\newcommand*\circled[1]{\tikz[baseline=(char.base)]{\node[shape=circle,draw,inner sep=1.5pt] (char) {\small #1};}}
\setcounter{MaxMatrixCols}{256}
\begin{document}
\section*{浦东新区 2025 学年度第一学期期末教学质量检测 高三数学试卷}

考生注意:1、本试卷共 21 道试题，满分 150 分，答题时间120 分钟；

2、请在答题纸上规定的地方解答，否则一律不予评分.

\section*{一、填空题(本大题满分 54 分)本大题共有 12 题.考生应在答题纸相应编号的 空格内直接填写结果, 1-6 题每个空格填对得 4 分, 7-12 题每个空格填对得 5 分, 否则一律得零分.}

1. 已知集合 \(A = \{ x \mid  x > 1\} ,B = \{ 1,2,3\}\) ，则 \(A \cap  B =\) \_\_\_\_\_.

【答案】 \(\{ 2,3\}\)

【解析】

【分析】

根据集合交集的定义和运算, 即可求解.

【详解】由题意,集合 \(A = \{ x \mid  x > 1\} ,B = \{ 1,2,3\}\) ,

根据集合交集的定义和运算,可得 \(A \cap  B = \{ 2,3\}\) .

故答案为: \(\{ 2,3\}\) .

【点睛】本题主要考查了集合交集的定义及运算, 其中熟记集合交集的定义是解答的关键, 属于容易题.

2. 函数 \(y = \lg \left( {x + 1}\right)\) 的定义域为\_\_\_\_\_

【答案】 \(\left( {-1, + \infty }\right)\)

【解析】

【分析】根据对数函数的定义域列不等式求解即可.

【详解】由 \(x + 1 > 0\) 得 \(x >  - 1\) ,

所以函数 \(y = \lg \left( {x + 1}\right)\) 的定义域为 \(\left( {-1, + \infty }\right)\) .

故答案为: \(\left( {-1, + \infty }\right)\) .

3. 已知向量 \(\overrightarrow{a} = \left( {1, - 3}\right)\) 、 \(\overrightarrow{b} = \left( {1,k}\right)\) ，若 \(\overrightarrow{a}//\overrightarrow{b}\) ，则实数 \(k\) 的值为\_\_\_\_\_.

【答案】 -3

【解析】

【分析】根据平面向量共线求参数即可.

【详解】因为 \(\overrightarrow{a}//\overrightarrow{b}\) ,所以 \(1 \cdot  k = \left( {-3}\right)  \times  1 \Rightarrow  k =  - 3\) .

故答案为: -3

4. 在 \({\left( x + \frac{1}{x}\right) }^{6}\) 的二项展开式中,常数项为\_\_\_\_\_.

【答案】 20

【解析】

【分析】根据二项式展开式的通项特征即可求解.

【详解】 \({\left( x + \frac{1}{x}\right) }^{6}\) 的二项展开式中,常数项为 \({\mathrm{C}}_{6}^{3}{x}^{3}{\left( \frac{1}{x}\right) }^{3} = {20}\) ,

故答案为: 20

5. 已知正数 \(a,b\) 满足 \(a + {2b} = 2\) ，则 \({ab}\) 的最大值为\_\_\_\_\_.

【答案】 \(\frac{1}{2}\# \# {0.5}\)

【解析】

【分析】直接利用基本不等式求解即可.

【详解】因为正数 \(a,b\) 满足 \(a + {2b} = 2\) ,所以 \(2 = a + {2b} \geq  2\sqrt{a \times  {2b}}\) ,平方化简得 \({ab} \leq  \frac{1}{2}\) , 当且仅当 \(a = {2b} = 1\) 时,等号成立,所以 \({ab}\) 的最大值为 \(\frac{1}{2}\) .

故答案为: \(\frac{1}{2}\)

6. 设 \(x \in  \mathbf{R}\) ,求方程 \(\left| {x + 2}\right|  + \left| {x - 2}\right|  = 4\) 的解集为\_\_\_\_\_.

【答案】[-2,2]

【解析】

【分析】根据绝对值三角不等式求解 \(\left| {x + 2}\right|  + \left| {x - 2}\right|\) 的最值,根据取等条件即可得方程的解集.

【详解】因为 \(\left| {x + 2}\right|  + \left| {2 - x}\right|  \geq  \left| {x + 2 + 2 - x}\right|  = 4\) ,

当且仅当 \(\left( {x + 2}\right) \left( {2 - x}\right)  \geq  0\) 时取等,

解得 \(x \in  \left\lbrack  {-2,2}\right\rbrack\) ,

所以方程 \(\left| {x + 2}\right|  + \left| {x - 2}\right|  = 4\) 的解集为 \(\left\lbrack  {-2,2}\right\rbrack\) .

故答案为: \(\left\lbrack  {-2,2}\right\rbrack\) .

7. \(\bigtriangleup  {ABC}\) 中， \(A = \frac{2\pi }{3},b = 1,\sin C = 2\sin B\) ，则 \(a =\) \_\_\_\_\_.

【答案】 \(\sqrt{7}\)

【解析】

【分析】根据正弦定理以及余弦定理, 可得答案.

【详解】由题意可得 \(\sin C = 2\sin B \Rightarrow  c = {2b} \Rightarrow  c = 2\) ,

因为 \(A = \frac{2\pi }{3}\) ,所以 \(\cos A = \frac{{b}^{2} + {c}^{2} - {a}^{2}}{2bc} \Rightarrow   - \frac{1}{2} = \frac{{1}^{2} + {2}^{2} - {a}^{2}}{2 \times  1 \times  2} \Rightarrow  a = \sqrt{7}\) .

故答案为: \(\sqrt{7}\) .

8. 如图,等腰直角 \(\bigtriangleup  {APB}\) 的斜边 \({AB}\) 长为 \({2R}\) ，将 \(\bigtriangleup  {APB}\) 绕斜边 \({AB}\) 所在直线旋转一周形成的旋转体的体积为\_\_\_\_\_.

\begin{center}
\includegraphics[max width=0.2\textwidth]{images/bo_d6dbjts601uc73dvo8s0_2_240_1183_239_368_0.jpg}
\end{center}
\hspace*{3em} 

【答案】 \(\frac{2}{3}\pi {R}^{3}\)

【解析】

【分析】旋转体为两个同底等高的圆锥的组合体, 利用体积公式求解即可

【详解】由题意,将 \(\bigtriangleup {APB}\) 绕斜边 \({AB}\) 所在直线旋转一周形成的旋转体是以 \(R\) 为底面半径, \(R\) 为高的两个共底面圆锥,

所以旋转体的体积为 \(V = 2 \times  \frac{1}{3}\pi {R}^{2} \cdot  R = \frac{2}{3}\pi {R}^{3}\) .

故答案为: \(\frac{2}{3}\pi {R}^{3}\)

9. 已知奇函数 \(y = f\left( x\right)\) 的定义域为 \(\mathbf{R}\) ,当 \(x < 0\) 时, \(f\left( x\right)  = {x}^{3} + {\mathrm{e}}^{x} - 1\) ,则函数 \(y = f\left( x\right)\) 在 \(x = 1\) 处的导数 \({f}^{\prime }\left( 1\right)  =\) \_\_\_\_\_.

【答案】 \(3 + {\mathrm{e}}^{-1}\)

【解析】

【分析】先根据奇函数的定义求出 \(x > 0\) 时的函数的解析式,再求导即可求出导数值.

【详解】因为奇函数 \(y = f\left( x\right)\) 的定义域为 \(\mathbf{R}\) ,

令 \(x > 0\) ,则 \(- x < 0\) ,所以 \(f\left( x\right)  =  - f\left( {-x}\right)  =  - \left( {-{x}^{3} + {\mathrm{e}}^{-x} - 1}\right)  = {x}^{3} - {\mathrm{e}}^{-x} + 1\) ,

所以当 \(x > 0\) 时, \(f\left( x\right)  = {x}^{3} - {\mathrm{e}}^{-x} + 1\) ,

所以 \({f}^{\prime }\left( x\right)  = 3{x}^{2} + {\mathrm{e}}^{-x} \Rightarrow  {f}^{\prime }\left( 1\right)  = 3 + {\mathrm{e}}^{-1}\) .

故答案为: \(3 + {\mathrm{e}}^{-1}\) .

10. 已知复数 \({z}_{1},{z}_{2}\) 是实系数一元二次方程的两个根,若 \({z}_{1} = \mathrm{i}{z}_{2}\) ,则 \(\left| {{z}_{1} - \mathrm{i}}\right|\) 的最小值为\_\_\_\_\_. 【答案】 \(\frac{\sqrt{2}}{2}\)

【解析】

【分析】分一元二次方程的判别式大于等于 0 与小于 0 , 两种情况讨论, 利用实系数一元二次方程的虚根成对的性质,计算可求得 \(\left| {{z}_{1} - \mathrm{i}}\right|\) 的最小值.

【详解】若一元二次方程的判别式大于等于 0,则方程有两个实数根,即 \({z}_{1},{z}_{2}\) 为实数, 由 \({z}_{1} = \mathrm{i}{z}_{2}\) ,则 \({z}_{1} = 0,{z}_{2} = 0\) ,此时 \(\left| {{z}_{1} - \mathrm{i}}\right|  = \left| {0 - \mathrm{i}}\right|  = 1\) ,

若一元二次方程的判别式小于 0,则 \({z}_{1},{z}_{2}\) 为两虚数根,

设 \({z}_{1} = a + b\mathrm{i}\text{ 、 }{z}_{2} = a - b\mathrm{i}\)

又因为 \({z}_{1} = \mathrm{i}{z}_{2}\) ,所以 \(a + b\mathrm{i} = \mathrm{i}\left( {a - b\mathrm{i}}\right)  = b + a\mathrm{i}\) ,所以 \(a = b\) ,

所以 \(\left| {{z}_{1} - \mathrm{i}}\right|  = \left| {a + \left( {a - 1}\right) \mathrm{i}}\right|  = \sqrt{{a}^{2} + {\left( a - 1\right) }^{2}} = \sqrt{2{a}^{2} - {2a} + 1} = \sqrt{2{\left( a - \frac{1}{2}\right) }^{2} + \frac{1}{2}}\)

当 \(a = \frac{1}{2}\) 时, \({\left| {z}_{1} - \mathrm{i}\right| }_{\min } = \sqrt{2{\left( \frac{1}{2} - \frac{1}{2}\right) }^{2} + \frac{1}{2}} = \frac{\sqrt{2}}{2}\) .

综上所述: \(\left| {{z}_{1} - \mathrm{i}}\right|\) 的最小值为 \(\frac{\sqrt{2}}{2}\) .

故答案为: \(\frac{\sqrt{2}}{2}\) .

11. 某城市交通系统采用数智技术记录城市道路口的信号灯状态, 每个时段的信号灯状态对应一个数字 (见下表).

\begin{center}
\adjustbox{max width=\textwidth}{
\begin{tabular}{|c|c|c|c|}
\hline
\multirow{2}{*}{信号灯状态} & \multirow{2}{*}{正常绿波协调} & \multicolumn{2}{|c|}{非正常绿波协调} \\
\cline{3-4}
 &  & 道路施工临时黄灯 & 突发事故临时红灯 \\
\cline{1-4}
对应数字 & 0 & 1 & -1 \\
\cline{1-4}
\hline
\end{tabular}
}
\end{center}

现需要按顺序记录某个道路口 5 个时段的信号灯状态,例如,记号 \(\left( {1,0, - 1,0,1}\right)\) 表示 5 个时段中有 3 个时段是“非正常绿波协调”状态，分别发生在第 1、3、5 时段.间:该路口 5 个时段所有可能的记录中，“非正常绿波协调”的时段数不少于 1 个且不多于 3 个的记录共有\_\_\_\_\_ 种.

【答案】130

【解析】

【分析】按照“非正常绿波协调”的时段个数进行分类, 结合计数原理和组合数知识求解.

【详解】按照“非正常绿波协调”的时段个数进行分类,

每类情况均先从 5 个时段中选取 “非正常绿波协调”的时段个数，且因 “非正常绿波协调”有两种情况，

则“非正常绿波协调”的时段数 1 个的记录有 \(2{\mathrm{C}}_{5}^{1} = {10}\) 种；

“非正常绿波协调”的时段数 2 个的记录有 \({2}^{2}{\mathrm{C}}_{5}^{2} = {40}\) 种；

“非正常绿波协调”的时段数 3 个的记录有 \({2}^{3}{\mathrm{C}}_{5}^{3} = {80}\) 种;

故 “非正常绿波协调”的时段数不少于 1 个且不多于 3 个的记录共有: \({10} + {40} + {80} = {130}\) 种; 故答案为: 130

12. 已知 \(f\left( x\right)  = \left\{  \begin{array}{l} 8{x}^{2} + {3x} - 1,x < \frac{1}{4} \\  x + \frac{7}{8},x \geq  \frac{1}{4} \end{array}\right.\) ,若等差数列 \(\left\{  {a}_{n}\right\}\) 为无穷数列,且均满足递推关系 \({a}_{n + 1} = f\left( {a}_{n}\right)\) ，则该数列首项 \({a}_{1}\) 的取值范围为\_\_\_\_\_.

【答案】 \(\left\{  {-\frac{1}{2}, - \frac{5}{8}}\right\}   \cup  \left\lbrack  {\frac{1}{4}, + \infty }\right)\) .

【解析】

【分析】由递推关系 \({a}_{n + 1} = f\left( {a}_{n}\right)\) ,则 \(f\left( {a}_{n}\right)  = {a}_{n + 1} = {a}_{n} + d\) ,根据分段函数的解析式,分三种情况讨论,列出关系式,即可求解.

【详解】由函数 \(f\left( x\right)  = \left\{  \begin{array}{l} 8{x}^{2} + {3x} - 1,x < \frac{1}{4} \\  x + \frac{7}{8},x \geq  \frac{1}{4} \end{array}\right.\) ,可得函数 \(f\left( x\right)\) 的图像,如图所示,

因为等差数列 \(\left\{  {a}_{n}\right\}\) 为无穷数列,且均满足递推关系 \({a}_{n + 1} = f\left( {a}_{n}\right)\) ,则 \(f\left( {a}_{n}\right)  = {a}_{n + 1} = {a}_{n} + d,\)

当 \({a}_{n} \in  \left\lbrack  {\frac{1}{4}, + \infty }\right)\) 时, \(f\left( {a}_{n}\right)  = {a}_{n + 1} = {a}_{n} + \frac{7}{8} = {a}_{n} + d\) ,

可得 \(d = \frac{7}{8}\) ,符合题意,此时 \({a}_{1} \in  \left\lbrack  {\frac{1}{4}, + \infty }\right)\) ;

当 \({a}_{n} \in  \left( {-\infty ,\frac{1}{4}}\right)\) 时,则 \({a}_{n + 1} \in  \left( {-\infty ,\frac{1}{4}}\right)\) ,可得 \(f\left( {a}_{n}\right)  = 8{a}_{n}^{2} + 3{a}_{n} - 1 = {a}_{n} + d\)

\(f\left( {a}_{n + 1}\right)  = 8{a}_{n + 1}^{2} + 3{a}_{n + 1} - 1 = {a}_{n + 1} + d,\)

两式相减,可得 \(8\left( {{a}_{n + 1}^{2} - {a}_{n}^{2}}\right)  + 2\left( {{a}_{n + 1} - {a}_{n}}\right)  = 0\) ,即 \(\left( {{a}_{n + 1} - {a}_{n}}\right) \left\lbrack  {4\left( {{a}_{n + 1} + {a}_{n}}\right)  + 1}\right\rbrack   = 0\) ,

所以 \({a}_{n + 1} - {a}_{n} = 0\) ,所以 \(d = 0\) ,此时数列 \(\left\{  {a}_{n}\right\}\) 为常数列,

可得 \(8{a}_{1}^{2} + 3{a}_{1} - 1 = {a}_{1}\) ,解得 \({a}_{1} =  - \frac{1}{2}\) 或 \({a}_{1} = \frac{1}{4}\) (舍去);

若 \({a}_{1} \in  \left( {-\infty ,\frac{1}{4}}\right) ,{a}_{2} \in  \left\lbrack  {\frac{1}{4}, + \infty }\right)\) ,则 \(f\left( {a}_{1}\right)  = 8{a}_{1}^{2} + 3{a}_{1} - 1 = {a}_{1} + \frac{7}{8}\) ,解得 \({a}_{1} =  - \frac{5}{8}\) 综上可得: 首项 \({a}_{1}\) 的取值范围为 \(\left\{  {-\frac{1}{2}, - \frac{5}{8}}\right\}   \cup  \left\lbrack  {\frac{1}{4}, + \infty }\right)\) .

故答案为: \(\left\{  {-\frac{1}{2}, - \frac{5}{8}}\right\}   \cup  \left\lbrack  {\frac{1}{4}, + \infty }\right)\) .

\begin{center}
\includegraphics[max width=0.2\textwidth]{images/bo_d6dbjts601uc73dvo8s0_6_240_193_329_310_0.jpg}
\end{center}
\hspace*{3em} 

\section*{二、选择题 (本大题满分 18 分) 本大题共有 4 题, 每题有且只有一个正确答案. 考生必须在答题纸的相应编号上, 将代表答案的小方格涂黑, 13-14 题每题选对 得 4 分, 15-16 题每题选对得 5 分, 否则一律得零分.}

13. 已知直线 \(a\text{ 、 }b\) 和平面 \(\alpha \text{ 、 }\beta\) ,且 \(a \subseteq  \alpha \text{ 、 }b \subseteq  \beta\) ,则 “ \(a\) 与 \(b\) 相交” 是 “ \(\alpha\) 与 \(\beta\) 相交”的 ( )

A. 充分必要条件 B. 既不充分又不必要条件

C. 充分不必要条件 D. 必要不充分条件

【答案】C

【解析】

【分析】根据直线与平面的位置关系判断.

【详解】直线 \(a\text{ 、 }b\) 和平面 \(\alpha \text{ 、 }\beta\) ,且 \(a \subseteq  \alpha \text{ 、 }b \subseteq  \beta\) ,

若 \(a\) 与 \(b\) 相交,记 \(a \cap  b = O\) ,由于 \(a \subseteq  \alpha \text{ 、 }b \subseteq  \beta\) ,

所以 \(O \in  \alpha \text{ 、 }O \in  \beta\) ,则 \(\alpha\) 与 \(\beta\) 相交,

所以 “ \(a\) 与 \(b\) 相交” 能得到 \(\alpha\) 与 \(\beta\) 相交,

但 “ \(\alpha\) 与 \(\beta\) 相交”,则 \(a\) 与 \(b\) 可能异面,如图,

\begin{center}
\includegraphics[max width=0.3\textwidth]{images/bo_d6dbjts601uc73dvo8s0_6_240_1606_391_304_0.jpg}
\end{center}
\hspace*{3em} 

平面 \({ABCD}\) 为 \(\alpha\) ,平面 \({ABEF}\) 为 \(\beta ,\alpha  \cap  \beta  = {AB}\) ,

\(a \subseteq  \alpha \text{ 、 }b \subseteq  \beta ,a \cap  {AB} = P,b \cap  {AB} = Q,\)

若 \(P\) 与 \(Q\) 不重合,则直线 \(a\text{ 、 }b\) 为异面直线,

所以 “ \(a\) 与 \(b\) 相交” 是 “ \(\alpha\) 与 \(\beta\) 相交” 的充分不必要条件.

故选: C.

14. 某班一次数学小测验 (百分制) 后, 老师为了奖励同学们平时认真学习, 决定给每位同学的成绩加上 5 分作为过程性评价奖励. 加分后, 与原始分数相比, 不会发生改变的是 ( )

A. 平均数 B. 中位数 C. 第 80 百分位数 D. 方差

【答案】D

【解析】

【分析】根据平均值、中位数、百分位数的概念判断 \(\mathrm{{ABC}}\) ,根据方差性质判断 \(\mathrm{D}\) .

【详解】加分后, 与原始分数相比, 平均值, 中位数, 第 80 百分位数的数值都会发生改变, 但根据方差的性质, 一组数据同时加上相同的数后, 方差大小不变.

故选: D.

15. 如图,椭圆 \({C}_{1}\text{ 、 }{C}_{2}\) 的离心率分别为 \({e}_{1}\text{ 、 }{e}_{2}\) ,双曲线 \({C}_{3}\text{ 、 }{C}_{4}\) 的离心率分别为 \({e}_{3}\text{ 、 }{e}_{4}\) , 则下列结论正确的是( )

\begin{center}
\includegraphics[max width=0.3\textwidth]{images/bo_d6dbjts601uc73dvo8s0_7_242_1090_374_295_0.jpg}
\end{center}
\hspace*{3em} 

A. \({e}_{3} > {e}_{4} > {e}_{2} > {e}_{1}\) B. \({e}_{4} > {e}_{3} > {e}_{1} > {e}_{2}\) C. \({e}_{3} > {e}_{4} > {e}_{1} > {e}_{2}\) D.

\({e}_{4} > {e}_{3} > {e}_{2} > {e}_{1}\)

【答案】D

【解析】

【分析】根据双曲线离心率: \(e = \frac{c}{a} = \sqrt{1 + \frac{{b}^{2}}{{a}^{2}}}\) ,椭圆离心率 \(e = \frac{c}{a} = \sqrt{1 - \frac{{b}^{2}}{{a}^{2}}}\) ,比较大小即可判断.

【详解】双曲线离心率: \(e = \frac{c}{a} = \sqrt{1 + \frac{{b}^{2}}{{a}^{2}}}\) ,椭圆离心率 \(e = \frac{c}{a} = \sqrt{1 - \frac{{b}^{2}}{{a}^{2}}}\) ,

如图可知椭圆 \({C}_{1},{C}_{2}\) 中 \({a}_{1} = {a}_{2},{b}_{2} < {b}_{1}\) ,所以 \({e}_{1} < {e}_{2} < 1\) ,

双曲线 \({C}_{3}\text{ 、 }{C}_{4}\) 中 \({a}_{3} > {a}_{4},{b}_{3} < {b}_{4}\) ,所以 \(1 < {e}_{3} < {e}_{4}\) ,

综上: \({e}_{4} > {e}_{3} > {e}_{2} > {e}_{1}\) .

故选: D.

16. 已知函数 \(y = f\left( x\right)\) ,其中 \(f\left( x\right)  = x - \frac{1}{x},x > 0\) . 集合 \(A = \{ 1,a,b\}\) ,其中 \(1 < a < b\) ,集合 \(B = \left\{  {z\left| {\;z = {xy} + \frac{x}{y} + f\left( \frac{y}{x}\right) }\right. ,x,y \in  A,x \neq  y}\right\}\) ,现有以下两个命题:

 \circled{1}  存在实数 \(a\text{ 、 }b\) ,使得集合 \(B\) 中恰好有 5 个元素;

 \circled{2}  若实数 \(a\text{ 、 }b\) 是方程 \(f\left( x\right)  = {kx} + t\) 的两个不同实根，则存在实数 \(k\text{ 、 }t\) ，使得集合 \(B\) 中恰好有 4 个元素.

那么 ( )

A.  \circled{1} 是真命题， \circled{2} 是真命题 B.  \circled{1} 是假命题， \circled{2} 是假命题

C.  \circled{1} 是真命题， \circled{2} 是假命题 D.  \circled{1} 是假命题， \circled{2} 是真命题

【答案】A

【解析】

【分析】先化简集合 \(B\) 的元素,通过构造具体数值验证命题 \circled{1} 中元素的合并情况; 结合韦达定理与函数零点存在性, 验证命题 \circled{2} 中元素的合并可行性。

【详解】 \(B = \left\{  {z\left| {\;z = {xy} + \frac{x}{y} + f\left( \frac{y}{x}\right) }\right. ,x,y \in  A,x \neq  y}\right\}   = \left\{  {z\left| {\;z = {xy} + \frac{y}{x}}\right. ,x,y \in  A,x \neq  y}\right\}\) , 集合 \(A = \{ 1,a,b\}\) ,其中 \(1 < a < b,x,y\) 取不同元素时,

集合 \(B\) 的元素可能为 \({2a},{2b},a + \frac{1}{a},b + \frac{1}{b},{ab} + \frac{b}{a},{ab} + \frac{a}{b}\) ,共 6 个可能.

 \circled{1} 若集合 \(B\) 中恰好有 5 个元素，则当 \({2a} = b + \frac{1}{b}\) 时， \(1 < a < b\) 必然能找到 \(a,b\) ，

例如: 当 \(a = \frac{5}{4}\) ,可取 \(b = 2\) ,此时 \(B = \left\{  {\frac{5}{2},4,\frac{41}{20},\frac{41}{10},\frac{25}{8}}\right\}\) ,共 5 个元素.

所以 \circled{1} 是真命题.

 \circled{2} 依题意，若存在实数 \(a\text{ 、 }b\) 是方程 \(f\left( x\right)  = {kx} + t\) 的两个不同实根，

\(x - \frac{1}{x} = {kx} + t \Rightarrow  \left( {k - 1}\right) {x}^{2} + {tx} + 1 = 0\)

所以 \(a + b = \frac{-t}{k - 1},{ab} = \frac{1}{k - 1}\) .

因为 \(1 < a < b,a + \frac{1}{a} < b + \frac{1}{b} < {2b},a + \frac{1}{a} < a\left( {b + \frac{1}{b}}\right) ,b + \frac{1}{b} < a\left( {b + \frac{1}{b}}\right)\)

若 \(\left\{  {\begin{array}{l} {2a} = b + \frac{1}{b} \\  {2b} = {ab} + \frac{b}{a} \end{array} \Rightarrow  \left\{  \begin{array}{l} a = 1 \\  b = 1 \end{array}\right. }\right.\) 不符;

若 \(\left\{  {\begin{array}{l} {2a} = b + \frac{1}{b} \\  {2b} = {ab} + \frac{a}{b} \end{array} \Rightarrow  {a}^{4} - 2{a}^{3} + 1 = 0}\right. \left( *\right)\) ,

设 \(h\left( x\right)  = {x}^{4} - 2{x}^{3} + 1,x \in  \left( {1, + \infty }\right)\)

\(h\left( \frac{3}{2}\right)  =  - \frac{11}{16},h\left( 2\right)  = 1\)

\(h\left( \frac{3}{2}\right)  \times  h\left( 2\right)  < 0\)

所以 \(h\left( x\right)  = {x}^{4} - 2{x}^{3} + 1,x \in  \left( {1, + \infty }\right)\) 有零点,即 (*) 有解,

即存在 \(a,b\) ,满足 \(\left\{  \begin{array}{l} {2a} = b + \frac{1}{b} \\  {2b} = {ab} + \frac{a}{b} \end{array}\right.\) ,使得 \(B\) 中恰好有 4 个元素.

所以实数 \(k\text{ 、 }t\) ,使得集合 \(B\) 中恰好有 4 个元素, \circled{2}  是真命题.

故选: A.

\section*{三、解答题(本大题满分 78 分)本大题共有 5 题，解答下列各题必须在答题卷 的相应编号规定区域内写出必要的步骤.}

17. 已知函数 \(y = f\left( x\right)\) ,其中 \(f\left( x\right)  = \cos \left( {{2x} - \frac{\pi }{3}}\right)\) .

(1)求函数 \(y = f\left( x\right)\) 的最小正周期及单调减区间.

(2)求函数 \(y = 2{\sin }^{2}x + f\left( x\right) ,x \in  \left\lbrack  {-\frac{\pi }{6},\frac{\pi }{3}}\right\rbrack\) 的最大值，并求出取得最大值时 \(x\) 的值.

【答案】(1) \(T = \pi ,\left\lbrack  {{k\pi } + \frac{\pi }{6},{k\pi } + \frac{2\pi }{3}}\right\rbrack  \left( {k \in  \mathbf{Z}}\right)\) ;

(2) \(2,x = \frac{\pi }{3}\) .

【解析】

【分析】(1)利用余弦函数的周期公式及单调性列式求解.

(2)利用三角函数恒等变换化简函数，再利用正弦函数的性质求出指定区间上的最大值.

【小问 1 详解】

函数 \(y = f\left( x\right)\) 的最小正周期 \(T = \frac{2\pi }{2} = \pi\) ,

由 \({2k\pi } \leq  {2x} - \frac{\pi }{3} \leq  {2k\pi } + \pi ,k \in  \mathbf{Z}\) ,得 \({k\pi } + \frac{\pi }{6} \leq  x \leq  {k\pi } + \frac{2\pi }{3},k \in  \mathbf{Z}\) ,

所以函数 \(f\left( x\right)\) 单调减区间为 \(\left\lbrack  {{k\pi } + \frac{\pi }{6},{k\pi } + \frac{2\pi }{3}}\right\rbrack  \left( {k \in  \mathbf{Z}}\right)\) .

【小问 2 详解】

依题意, \(y = 2{\sin }^{2}x + \cos \left( {{2x} - \frac{\pi }{3}}\right)  = 1 - \cos {2x} + \cos {2x}\cos \frac{\pi }{3} + \sin {2x}\sin \frac{\pi }{3}\)

所以 \(= \frac{\sqrt{3}}{2}\sin {2x} - \frac{1}{2}\cos {2x} + 1 = \sin \left( {{2x} - \frac{\pi }{6}}\right)  + 1\) ,

由 \(x \in  \left\lbrack  {-\frac{\pi }{6},\frac{\pi }{3}}\right\rbrack\) ,得 \({2x} - \frac{\pi }{6} \in  \left\lbrack  {-\frac{\pi }{2},\frac{\pi }{2}}\right\rbrack\) ,则当 \({2x} - \frac{\pi }{6} = \frac{\pi }{2}\) ,即 \(x = \frac{\pi }{3}\) 时,函数取得最大值 2, 所以最大值为 2,此时 \(x = \frac{\pi }{3}\) .

18. 如图,该几何体由一个棱长为 4 的正方体 \({ABCD} - {A}_{1}{B}_{1}{C}_{1}{D}_{1}\) 与一个半圆柱拼接而成, 圆心 \(O\text{ 、 }{O}_{1}\) 分别为线段 \({AB}\text{ 、 }{A}_{1}{B}_{1}\) 的中点,动点 \(M\) 在弧 \(\overset{\text{ ⏜ }}{{A}_{1}{B}_{1}}\) 上滑动.

\begin{center}
\includegraphics[max width=0.3\textwidth]{images/bo_d6dbjts601uc73dvo8s0_10_240_1475_402_349_0.jpg}
\end{center}
\hspace*{3em} 

(1)若点 \(M\) 为弧 \(\overset{\text{ ⏜ }}{{A}_{1}{B}_{1}}\) 的中点,求直线 \({MD}\) 与平面 \({BD}{D}_{1}{B}_{1}\) 所成角的大小;

(2)若弧 \(\overset{\text{ ⏜ }}{{A}_{1}M}\) 的长度为 \(\frac{\pi }{2}\) ，求证:平面 \(O{O}_{1}M//\) 平面 \({AC}{C}_{1}{A}_{1}\) .

【答案】(1) \(\arcsin \frac{\sqrt{7}}{7}\)

(2)证明见解析

【解析】

【分析】(1) 建立空间直角坐标系,求得 \(\overrightarrow{MD}\) 与平面 \({BD}{D}_{1}{B}_{1}\) 的法向量,通过向量夹角公式求线面角;

(2)由弧长可求得 \(M\) 点坐标，进而求得平面 \(O{O}_{1}M\) 和平面 \({AC}{C}_{1}{A}_{1}\) 的法向量，计算即可得到结果.

【小问 1 详解】

建立如图所示空间直角坐标系:

则: \(A\left( {0,0,0}\right) ,B\left( {4,0,0}\right) ,D\left( {0,4,0}\right) ,C\left( {4,4,0}\right) ,O\left( {2,0,0}\right) ,{O}_{1}\left( {2,0,4}\right)\)

\begin{center}
\includegraphics[max width=0.3\textwidth]{images/bo_d6dbjts601uc73dvo8s0_11_238_863_422_405_0.jpg}
\end{center}
\hspace*{3em} 

若点 \(M\) 为弧 \(\overset{\text{ ⏜ }}{{A}_{1}{B}_{1}}\) 的中点,则 \(M\left( {2, - 2,4}\right)\)

所以 \(\overrightarrow{MD} = \left( {-2,6, - 4}\right)\)

平面 \({BD}{D}_{1}{B}_{1}\) 的一个法向量为 \(\overrightarrow{m} = \overrightarrow{AC} = \left( {4,4,0}\right)\)

所以 \(\sin \theta  = \cos \overrightarrow{m},\overrightarrow{MD} = \left| \frac{\overrightarrow{m} \cdot  \overrightarrow{MD}}{\left| \overrightarrow{m}\right|  \cdot  \left| \overrightarrow{MD}\right| }\right|  = \left| \frac{16}{4\sqrt{2} \times  2\sqrt{14}}\right|  = \frac{\sqrt{7}}{7}\)

所以直线 \({MD}\) 与平面 \({BD}{D}_{1}{B}_{1}\) 所成角的大小 \(\arcsin \frac{\sqrt{7}}{7}\)

【小问 2 详解】

若弧 \(\overset{\text{ ⏜ }}{{A}_{1}M}\) 的长度为 \(\frac{\pi }{2}\) ,由弧长得 \(\angle M{O}_{1}{A}_{1} = \frac{\frac{\pi }{2}}{2} = \frac{\pi }{4}\) ,则 \(M\left( {2 - \sqrt{2}, - \sqrt{2},4}\right)\)

\(\overrightarrow{O{O}_{1}} = \left( {0,0,4}\right) ,\overrightarrow{OM} = \left( {-\sqrt{2}, - \sqrt{2},4}\right)\)

设平面 \(O{O}_{1}M\) 的一个法向量为 \({\overrightarrow{n}}_{1} = \left( {x,y,z}\right)\) ,平面 \({AC}{C}_{1}{A}_{1}\) 的一个法向量为

\[
{\overrightarrow{n}}_{2} = \overrightarrow{BD} = \left( {-4,4,0}\right)
\]

\[
\left\{  {\begin{array}{l} {\overrightarrow{n}}_{1} \cdot  \overrightarrow{O{O}_{1}} = 0 \\  {\overrightarrow{n}}_{1} \cdot  \overrightarrow{OM} = 0 \end{array} \Rightarrow  \left\{  \begin{matrix} z = 0 \\  x + y = 0 \end{matrix}\right. }\right. \text{ 令 }x =  - 1
\]

\({\overrightarrow{n}}_{1} = \left( {-1,1,0}\right)\)

所以 \({\overrightarrow{n}}_{1}//{\overrightarrow{n}}_{2}\)

所以平面 \(O{O}_{1}M//\) 平面 \({AC}{C}_{1}{A}_{1}\)

19. 小明和小李要进行一系列的比赛, 假设每局比赛的结果互不影响.

(1)若比赛没有平局，且小明每局获胜的概率为 0.4 .

 \circled{1} 如果共有三局比赛，求小明的比赛结果依次为赢、输、赢的概率;

 \circled{2} 小明作为实力较弱的一方，他可以优先选择“一局定胜负”或“三局两胜”的赛制(“三局两胜”指先赢两局者为胜, 最多三局结束). 请帮助小明分析, 选择哪种赛制对他更有利, 并说明理由.

(2)如果小明每局获胜的概率为 \(p\left( {0 < p < 1}\right)\) ，他和小李要进行一场“五局三胜”的比赛(“五局三胜”指先赢三局者为胜,最多五局结束). 记小明最终获胜的概率为 \(f\left( p\right)\) ,请给出 \(f\left( p\right)\) 的表达式,判断并说明函数 \(y = f\left( p\right)\) 在 \(\left( {0,1}\right)\) 上的单调性,并指出现实意义.

【答案】(1)  \circled{1}  \(\frac{12}{125}\) ;  \circled{2} 一局定胜负对小明更有利,理由见解析

(2) \(f\left( p\right)  = 6{p}^{5} - {15}{p}^{4} + {10}{p}^{3}\) ，函数 \(y = f\left( p\right)\) 在 \(\left( {0,1}\right)\) 上的单调递增；意义见解析

【解析】

【分析】(1)  \circled{1} 根据独立事件概率的乘法公式即可计算小明比赛结果依次为 “赢、输、赢” 的概率;

 \circled{2}  分别计算 “一局定胜负” 和 “三局两胜” 赛制下小明获胜的概率, 比较大小来确定哪种赛制对小明更有利;

(2)先根据 “五局三胜” 的规则，分情况讨论小明获胜的情况，利用独立事件概率的乘法公式和互斥事件概率的加法公式,求出 \(f\left( p\right)\) 的表达式,再通过求导判断函数的单调性,并分析现实意义.

【小问 1 详解】

 \circled{1} 根据题意，比赛没有平局，且小明每局获胜的概率为 0.4 ，所以三局比赛相互独立，且小明每局输的概率为 0.6 ,

所以小明的比赛结果依次为赢、输、赢的概率为 \(P = \frac{2}{5} \times  \frac{3}{5} \times  \frac{2}{5} = \frac{12}{125}\) ;

 \circled{2} 根据题意，一局定胜负时，小明赢得比赛的概率为 \(p = \frac{2}{5}\) ，

三局两胜时, 小明赢得比赛有两种情况:

情况一:前两局获胜，概率为 \(\frac{2}{5} \times  \frac{2}{5} = \frac{4}{25}\) ，

情况二: 前两局胜一局,输一局,第三局获胜,其概率为 \({\mathrm{C}}_{2}^{1}{\left( \frac{2}{5}\right) }^{2}\left( \frac{3}{5}\right)  = \frac{24}{125}\) ,

根据互斥事件的加法公式,所以在 “三局两胜” 的赛制下,

小明获胜的概率为 \(\frac{4}{25} + \frac{24}{125} = \frac{44}{125}\) ,

因为 \(\frac{2}{5} > \frac{44}{125}\) ,所以从小明的角度考虑,一局定胜负对小明更有利;

【小问 2 详解】

“五局三胜”的赛制下，小明获胜有以下几种情况:

情况一:前三局获胜，概率为 \({P}_{1} = {p}^{3}\) ，

情况二:前三局胜两局输一局，第四局获胜，则 \({P}_{2} = {\mathrm{C}}_{3}^{2}{p}^{3}\left( {1 - p}\right)\) ，

情况三: 前四局胜两局输两局,第五局获胜,则 \({P}_{3} = {\mathrm{C}}_{4}^{2}{p}^{3}{\left( 1 - p\right) }^{2}\) ,

所以小明赢的概率为:

\(P = {P}_{1} + {P}_{2} + {P}_{3} = {p}^{3} + 3{p}^{3}\left( {1 - p}\right)  + 6{p}^{3}{\left( 1 - p\right) }^{2} = {p}^{3}\left( {6{p}^{2} - {15p} + {10}}\right) ,\)

所以 \(f\left( p\right)  = 6{p}^{5} - {15}{p}^{4} + {10}{p}^{3}\) ,

所以 \({f}^{\prime }\left( p\right)  = {30}{p}^{4} - {60}{p}^{3} + {30}{p}^{2} = {30}{p}^{2}{\left( p - 1\right) }^{2}\) ,

因为 \(0 < p < 1\) ,所以 \({f}^{\prime }\left( p\right)  = {30}{p}^{2}{\left( p - 1\right) }^{2} > 0\) ,

所以函数 \(y = f\left( p\right)\) 在 \(\left( {0,1}\right)\) 上的单调递增;

现实意义:在 “五局三胜” 的比赛中，小明每局获胜的概率 \(p\) 越大，他最终获胜的概率 \(f\left( p\right)\) 就越大, 即小明实力越强, 他获胜的可能性就越大.

20. 已知椭圆 \(C : \frac{{x}^{2}}{9} + \frac{{y}^{2}}{8} = 1\) ，点 \(O\) 为坐标原点，椭圆的右顶点为 \(A\) ，左右焦点分别为 \({F}_{1}\) 、 \({F}_{2}\) ， 点 \(E\) 为椭圆在 \(x\) 轴上方的一动点.

(1)求椭圆 \(C\) 的焦距与 \({\bigtriangleup E}{F}_{1}{F}_{2}\) 的周长;

(2)若 \(\overrightarrow{EA} \cdot  \overrightarrow{E{F}_{1}} = \frac{9}{4}\) ，求点 \(E\) 到 \(x\) 轴的距离；

(3)过点 \({F}_{1}\) 作斜率为 \(k\) 的直线 \(l\) 交 \(y\) 轴正半轴于点 \(P\) ，点 \(P\) 位于椭圆内. 直线 \(l\) 与椭圆 \(C\) 交于 \(G\text{ 、 }H\) 两点,记 \(\bigtriangleup {OGH}\text{ 、 }\bigtriangleup O{F}_{1}P\) 的面积分别为 \({S}_{1}\text{ 、 }{S}_{2}\) ,求 \(\frac{{S}_{1}}{{S}_{2}}\) 的取值范围.

【答案】(1)焦距为2， \(\bigtriangleup E{F}_{1}{F}_{2}\) 的周长 8 ;

(2) \(\sqrt{6}\) ；

(3) \(\left( {\frac{9}{5},6}\right)\) .

【解析】

【分析】(1)求出椭圆 \(C\) 的长短半轴长及半焦距,再结合椭圆定义求解.

(2)由(1)求出点 \(A,{F}_{1}\) 坐标，设出点 \(E\) 坐标，再利用数量积的坐标表示列式求解.

(3)设直线 \(l\) 的方程为 \(x = {my} - 1\) ，求出 \(m\) 范围，联立直线与椭圆方程，结合韦达定理表示 出 \(\frac{{S}_{1}}{{S}_{2}}\) ,再利用对勾函数单调性求出范围.

【小问 1 详解】

椭圆 \(C : \frac{{x}^{2}}{9} + \frac{{y}^{2}}{8} = 1\) 的长短半轴长分别为 \(a = 3,b = 2\sqrt{2}\) ,则半焦距 \(c = \sqrt{{a}^{2} - {b}^{2}} = 1\) , 所以椭圆 \(C\) 的焦距 \({2c} = 2,\bigtriangleup E{F}_{1}{F}_{2}\) 的周长 \({2a} + {2c} = 8\) .

【小问 2 详解】

由 (1) 得 \(A\left( {3,0}\right) ,{F}_{1}\left( {-1,0}\right)\) ,设 \(E\left( {{x}_{0},{y}_{0}}\right) ,{y}_{0} > 0\) ,则 \(\frac{{x}_{0}^{2}}{9} + \frac{{y}_{0}^{2}}{8} = 1\) ,即 \({y}_{0}^{2} = 8 - \frac{8{x}_{0}^{2}}{9}\) ,

\(\overrightarrow{EA} = \left( {3 - {x}_{0}, - {y}_{0}}\right) ,\overrightarrow{E{F}_{1}} = \left( {-1 - {x}_{0}, - {y}_{0}}\right)\) ，由 \(\overrightarrow{EA} \cdot  \overrightarrow{E{F}_{1}} = \frac{9}{4}\) ，

得 \(\overrightarrow{EA} \cdot  \overrightarrow{E{F}_{1}} = \left( {3 - {x}_{0}}\right) \left( {-1 - {x}_{0}}\right)  + {y}_{0}^{2} = {x}_{0}^{2} - 2{x}_{0} - 3 + {y}_{0}^{2} = \frac{1}{9}{x}_{0}^{2} - 2{x}_{0} + 5 = \frac{9}{4}\) ,

整理得 \(4{x}_{0}^{2} - {72}{x}_{0} + {99} = 0\) ,而 \(- 3 < {x}_{0} < 3\) ,解得 \({x}_{0} = \frac{3}{2},{y}_{0} = \sqrt{6}\) ,点 \(E\left( {\frac{3}{2},\sqrt{6}}\right)\) , 所以点 \(E\) 到 \(x\) 轴的距离为 \(\sqrt{6}\) .

【小问 3 详解】

设直线 \(l\) 的方程为 \(x = {my} - 1\) ,因为直线 \(l\) 与 \(y\) 轴正半轴相交,故 \(m \neq  0\) ,

则 \(P\left( {0,\frac{1}{m}}\right)\) ,由点 \(P\) 在 \(y\) 轴正半轴上,且位于椭圆内,

得 \(0 < \frac{1}{m} < 2\sqrt{2}\) ,即 \(m > \frac{\sqrt{2}}{4}\) ,

由 \(\left\{  \begin{array}{l} x = {my} - 1 \\  8{x}^{2} + 9{y}^{2} = {72} \end{array}\right.\) ,消去 \(x\) 得 \(\left( {8{m}^{2} + 9}\right) {y}^{2} - {16my} - {64} = 0\) ,

\(\Delta  = {256}{m}^{2} + {256}\left( {8{m}^{2} + 9}\right)  > 0\;\) ,设 \(\;G\left( {{x}_{1},{y}_{1}}\right) ,H\left( {{x}_{2},{y}_{2}}\right) \;\) ,则

\({y}_{1} + {y}_{2} = \frac{16m}{8{m}^{2} + 9},{y}_{1}{y}_{2} = \frac{-{64}}{8{m}^{2} + 9},\)

\(\left| {{y}_{1} - {y}_{2}}\right|  = \sqrt{{\left( {y}_{1} + {y}_{2}\right) }^{2} - 4{y}_{1}{y}_{2}} = \sqrt{{\left( \frac{16m}{8{m}^{2} + 9}\right) }^{2} + \frac{256}{8{m}^{2} + 9}} = \frac{{48}\sqrt{{m}^{2} + 1}}{8{m}^{2} + 9}\) ,

\({S}_{1} = {S}_{\bigtriangleup {OGH}} = \frac{1}{2}\left| {O{F}_{1}}\right|  \cdot  \left| {{y}_{1} - {y}_{2}}\right|  = \frac{1}{2}\left| {{y}_{1} - {y}_{2}}\right| ,{S}_{2} = {S}_{\bigtriangleup O{F}_{1}P} = \frac{1}{2}\left| {O{F}_{1}}\right|  \cdot  \frac{1}{m} = \frac{1}{2m}\) ,

因此 \(\frac{{S}_{1}}{{S}_{2}} = m\left| {{y}_{1} - {y}_{2}}\right|  = \frac{{48m}\sqrt{{m}^{2} + 1}}{8{m}^{2} + 9} = {48}\sqrt{\frac{{m}^{2}\left( {{m}^{2} + 1}\right) }{{\left( 8{m}^{2} + 9\right) }^{2}}}\) ,令 \({m}^{2} = t > \frac{1}{8},\frac{{S}_{1}}{{S}_{2}} = {48}\sqrt{\frac{t\left( {t + 1}\right) }{{\left( 8t + 9\right) }^{2}}}\) ,

而 \(f\left( t\right)  = \frac{t\left( {t + 1}\right) }{{\left( 8t + 9\right) }^{2}} = \frac{1 + \frac{1}{t}}{{\left( 8 + \frac{9}{t}\right) }^{2}}\) ,令 \(1 + \frac{1}{t} = u \in  \left( {1,9}\right)\) ,则 \(f\left( t\right)  = \frac{u}{{\left( 9u - 1\right) }^{2}} = \frac{1}{{81u} + \frac{1}{u} - {18}}\) ,

函数 \(y = {81u} + \frac{1}{u} - {18}\) 在 \(\left( {1,9}\right)\) 上单调递增, \({64} < {81u} + \frac{1}{u} - {18} < \frac{6400}{9}\) ,

因此 \(\frac{9}{6400} < \frac{t\left( {t + 1}\right) }{{\left( 8t + 9\right) }^{2}} < \frac{1}{64},\frac{3}{80} < \sqrt{\frac{t\left( {t + 1}\right) }{{\left( 8t + 9\right) }^{2}}} < \frac{1}{8}\) ,则 \(\frac{9}{5} < \frac{{S}_{1}}{{S}_{2}} < 6\) ,

所以 \(\frac{{S}_{1}}{{S}_{2}}\) 的取值范围是 \(\left( {\frac{9}{5},6}\right)\) .

\begin{center}
\includegraphics[max width=0.4\textwidth]{images/bo_d6dbjts601uc73dvo8s0_16_236_195_468_438_0.jpg}
\end{center}
\hspace*{3em} 

21. 已知函数 \(y = f\left( x\right)\) 的定义域为 \(D\) . 若对任意区间 \(I \subseteq  D\) ,存在正整数 \(k\) ,使得 \({I}_{k} \cap  I \neq  \varnothing\) ,则称函数 \(y = f\left( x\right)\) 为“回归函数”, \(k\) 为回归指数. 其中 \({I}_{n} = \left\{  {{f}^{\left( n\right) }\left( x\right)  \mid  x \in  I}\right\}\) , \({f}^{\left( 1\right) }\left( x\right)  = f\left( x\right) ,{f}^{\left( 2\right) }\left( x\right)  = f\left( {f\left( x\right) }\right) ,\ldots \ldots ,{f}^{\left( n + 1\right) }\left( x\right)  = f\left( {{f}^{\left( n\right) }\left( x\right) }\right) .\)

(1)若 \(f\left( x\right)  = \frac{1}{x},g\left( x\right)  = {x}^{2}\) . 请分别判断函数 \(y = f\left( x\right) \text{ 、 }y = g\left( x\right)\) 是否为回归指数为 2 的“回归函数”;

(2)若 \(f\left( x\right)  = \frac{x + 1}{1 - x},D = \left\{  {x\left| {\;x \in  \mathbf{R}}\right. \text{ ,且 }x \neq  0,x \neq  1,x \neq   - 1}\right\}\) . 求证: 函数 \(y = f\left( x\right)\) 是“回归函数”，并求满足条件的回归指数 \(k\) 的最小值.

(3) 若 \(y = f\left( x\right)\) 是定义域为 \((0,1\rbrack\) 的函数,且 \(f\left( x\right)  = \left\{  \begin{array}{ll} \frac{x}{a}, & 0 < x \leq  a, \\  \frac{x - a}{1 - a}, & a < x \leq  1. \end{array}\right.\) ,其中 \(a \in  \left( {0,1}\right)\) . 请判断函数 \(y = f\left( x\right)\) 是否为 “回归函数”? 若是,请写出证明过程; 若不是,请说明理由.

【答案】(1) \(y = f\left( x\right)\) 是回归指数为 2 的 "回归函数 ", \(y = g\left( x\right)\) 不是回归指数为 2 的 " 回归函数 ";

(2)证明见解析， \(k\) 的最小值为 4；

(3)函数 \(y = f\left( x\right)\) 是 “回归函数”，证明见解析.

【解析】

【分析】(1) 计算出 \({f}^{\left( 1\right) }\left( x\right)  = \frac{1}{x},{f}^{\left( 2\right) }\left( x\right)  = x\) ,从而得到 \({I}_{2} \cap  I \neq  \varnothing\) 即可判断 \(y = f\left( x\right)\) , 再计算 \({g}^{\left( 2\right) }\left( x\right)  = {x}^{4}\) ,取 \(I = \left\lbrack  {2,3}\right\rbrack\) 即可判断;

(2)计算得 \({f}^{\left( 5\right) }\left( x\right)  = {f}^{\left( 1\right) }\left( x\right) ,\cdots ,{f}^{\left( k + 4\right) }\left( x\right)  = {f}^{\left( k\right) }\left( x\right)\) ,再取 \(I = \left( {2,3}\right)\) ,分析得 \({I}_{{4m} - 3} \cap  I = \varnothing ,{I}_{{4m} - 2} \cap  I = \varnothing ,{I}_{{4m} - 1} \cap  I = \varnothing\) ,则得到 \(k\) 的最小值;

(3)首先求出函数 \(y = f\left( x\right)\) 的值域为 \((0,1\rbrack\) ，再分区间求出对应直线斜率，再设区间 \(I\) 的长度为 \(L,L > 0\) ,根据变换则区间长度也会变化,最后分析得总存在这样的正整数 \(k\) ,使得 \({I}_{k} \cap  I \neq  \varnothing\) .

【小问 1 详解】

由题意, \(y = f\left( x\right)\) 的定义域 \(D = \left( {-\infty ,0}\right)  \cup  \left( {0, + \infty }\right)\) .

\({f}^{\left( 1\right) }\left( x\right)  = f\left( x\right)  = \frac{1}{x},{f}^{\left( 2\right) }\left( x\right)  = f\left( {f\left( x\right) }\right)  = f\left( \frac{1}{x}\right)  = x,\)

故对任意区间 \(I \subseteq  \left( {-\infty ,0}\right)  \cup  \left( {0, + \infty }\right)\) ,必有 \({I}_{2} = I\) ,即 \({I}_{2} \cap  I \neq  \varnothing\) .

因此 \(y = f\left( x\right)\) 是回归指数为 2 的 "回归函数 ".

又 \(y = g\left( x\right)\) 的定义域 \(D = \mathrm{R},{g}^{\left( 2\right) }\left( x\right)  = {\left( {x}^{2}\right) }^{2} = {x}^{4}\) .

取区间 \(I = \left\lbrack  {2,3}\right\rbrack\) ,则 \({I}_{2} = \left\lbrack  {{16},{81}}\right\rbrack\) . 显然, \(I \cap  {I}_{2} = \varnothing\) .

显然, \(y = g\left( x\right)\) 不是回归指数为 2 的 "回归函数 ".

【小问 2 详解】

由题意,对任意 \(x \in  D,{f}^{\left( 1\right) }\left( x\right)  = \frac{x + 1}{1 - x},{f}^{\left( 2\right) }\left( x\right)  = \frac{\frac{x + 1}{1 - x} + 1}{1 - \frac{x + 1}{1 - x}} = \frac{2}{-{2x}} =  - \frac{1}{x}\) ,

\({f}^{\left( 3\right) }\left( x\right)  = \frac{-\frac{1}{x} + 1}{1 - \left( {-\frac{1}{x}}\right) } = \frac{x - 1}{x + 1},{f}^{\left( 4\right) }\left( x\right)  = \frac{\frac{x - 1}{x + 1} + 1}{1 - \frac{x - 1}{x + 1}} = \frac{2x}{2} = x\) ,

\({f}^{\left( 5\right) }\left( x\right)  = {f}^{\left( 1\right) }\left( x\right) ,\cdots ,{f}^{\left( k + 4\right) }\left( x\right)  = {f}^{\left( k\right) }\left( x\right) .\)

显然,对任意区间 \(I \subseteq  D\) ,都存在正整数 \(m\) ,使得 \({I}_{4m} = I\) ,即 \({I}_{4m} \cap  I \neq  \varnothing\)

若取区间 \(I = \left( {2,3}\right)\) ,则 \({I}_{1} = \left( {-3, - 2}\right) ,{I}_{2} = \left( {\frac{1}{3},\frac{1}{2}}\right) ,{I}_{3} = \left( {-\frac{1}{2}, - \frac{1}{3}}\right)\) .

即 \({I}_{{4m} - 3} = {I}_{1} = \left( {-3, - 2}\right) ,{I}_{{4m} - 2} = {I}_{2} = \left( {\frac{1}{3},\frac{1}{2}}\right) ,{I}_{{4m} - 1} = {I}_{3} = \left( {-\frac{1}{2}, - \frac{1}{3}}\right)\) ,

故 \({I}_{{4m} - 3} \cap  I = \varnothing ,{I}_{{4m} - 2} \cap  I = \varnothing ,{I}_{{4m} - 1} \cap  I = \varnothing\) .

综上,函数 \(y = f\left( x\right)\) 是 “回归函数 ”,且满足条件的回归指数 \(k\) 的最小值为 4 .

【小问 3 详解】

当 \(x \in  (0,a\rbrack\) 时, \(f\left( x\right)  = \frac{x}{a},f\left( x\right)  \in  (0,1\rbrack\) .

当 \(x \in  (a,1\rbrack\) 时, \(f\left( x\right)  = \frac{x - a}{1 - a},f\left( x\right)  \in  (0,1\rbrack\) .

因此,函数 \(y = f\left( x\right)\) 的值域为 \((0,1\rbrack\) ,即对任意 \(y \in  (0,1\rbrack\) ,总存在 \(x \in  (0,1\rbrack\) 使得 \(f\left( x\right)  = y\) .

在区间 \((0,a\rbrack\) 上,函数为 \(f\left( x\right)  = \frac{x}{a}\) ,对应直线的斜率 \(k = \frac{1}{a} > 1\) .

在区间 \((a,1\rbrack\) 上, \(f\left( x\right)  = \frac{x - a}{1 - a}\) ,对应直线的斜率 \(k = \frac{1}{1 - a} > 1\) .

显然,对任意区间 \(I \subseteq  D\) ,作用函数 \(y = f\left( x\right)\) 后,区间长度必然会变大.

假设常数 \(c = \min \left\{  {\frac{1}{a},\frac{1}{1 - a}}\right\}\) ,必有 \(c > 1\) ,即每次变换后区间长度至少放大 \(c\) 倍.

设区间 \(I\) 的长度为 \(L,L > 0\) . 则经过 \(k\) 次变换后, \({I}_{k}\) 的长度至少为 \(L \cdot  {c}^{k}\) .

当 \(k\) 足够大时,必有 \(L \cdot  {c}^{k} > 1\) . 又函数 \(y = f\left( x\right)\) 的值域为 \((0,1\rbrack\) ,

即 \({I}_{k} \subseteq  (0,1\rbrack\) ,所以 \({I}_{k}\) 的长度不能超过 1,矛盾.

实际上,由于 \(y = f\left( x\right)\) 的图像是分段线性且连续的,变换过程中 \({I}_{k}\) 会覆盖 \((0,1\rbrack\) 的每个非空开区间子集.

必存在正整数 \(k\) ,使得 \({I}_{k} \cap  I \neq  \varnothing\) .

否则,如果所有 \({I}_{k}\) 都与 \(I\) 不相交,则 \({I}_{k}\) 始终停留在区间 \((0,1\rbrack\) 或区间 \(I\) 中,

但区间长度作用函数 \(y = f\left( x\right)\) 后会变大,即 \({I}_{k}\) 长度的不断增长,最终必须与区间 \(I\) 相交.

因此,总存在这样的正整数 \(k\) ,使得 \({I}_{k} \cap  I \neq  \varnothing\) .

综上所述,函数 \(y = f\left( x\right)\) 是 “回归函数”.
\end{document}